\honest{What's a Bias?}{Eliezer Yudkowsky}{2006}

I open with two examples of cognitive bias, availability and scope insensitivity. \nb{The text here is a later revision of the 2006 post, dated 2019 on LessWrong; the posts on availability and scope insensitivity were written in 2007.}

Cognitive biases are obstacles to truth produced ``not by the cost of information, nor by limited computing power, but by the shape of our own mental machinery.'' The machinery might have evolved to win arguments, or to believe what others believe, or a useful shortcut might have side effects. The availability heuristic is my example of the last: it does good work but also makes systematic errors. \nb{A shortcut is used because full calculation costs too much. The post does not say how its errors differ from errors caused by ``limited computing power,'' which it excludes.}

Once someone has identified such a problem ``verbally and concretely,'' we call it a bias. ``In cognitive science,'' biases are distinguished from mistakes, which come from learned false beliefs and are ``much easier to correct.'' \nb{No source is given for the claim about cognitive science, and no evidence that mistakes are easier to correct.} Biases are also distinguished from brain damage and absorbed culture: they come from machinery that is ``humanly universal.''

Plato was not biased because he did not know about relativity. But if he believed in philosopher-kings because of a universal instinct for self-promotion, and not because his father taught him or because he sniffed glue, then that was a bias. I do not say how anyone could find out which it was.

Having spent five paragraphs on the definition, I admit that the project does not rest on it: if ``cognitive bias'' turns out to be unhelpful, ``we should just drop it.''

We have no moral duty to reduce bias because biases are ``bad and evil and Just Not Done.'' Half in jest: doing techniques without knowing their reasons is itself ``bad and evil and Just Not Done, according to'' Feynman. A bias matters because it is in the way of the truth, which we need and are curious about.
